\section{Predictions and open problems}\label{sec:predictions}

This section records predictions and open problems.  None of them is
used as evidence for a numbered statement, and none asserts that a
missing application has been carried out.

\subsection{A Hodge substrate}
\label{sec:predictions:hodge}

We conjecture, on structural grounds only, that a Hodge substrate
\cite{Deligne2000Hodge,Voisin2003HodgeTheoryII} would lie in the
trace-state-only column; no formed closure for the Hodge conjecture has
been constructed in this series.  If one is constructed and the column
placement holds, the seven-stage template of \Cref{sec:recognition}
should apply to it, and its derivation sweep would be a third test of
the verdict signature \SigPattern\ of \Cref{thm:verdict-signature}.
The prediction concerns the column, not the substrate: it asserts no
Hodge closure and names no Hodge recognition source.  It would be
refuted by a trace-state-only Hodge closure satisfying the template
whose derivation sweep returns a different verdict pattern; if a Hodge
substrate turned out not to be trace-state-only, the prediction would
simply not apply.

\subsection{Yang--Mills}
\label{sec:predictions:ym}

Yang--Mills \cite{JaffeWitten2000YangMills} is a further candidate.
No Yang--Mills substrate has been studied in this series, and nothing
is claimed about it.  The open question is whether a Yang--Mills
substrate would be trace-state-only, witness-content-exposing, or of a
kind not yet represented.  The column question should be settled before
a recognition-mode or derivation-mode template is chosen.

\subsection{Proving the recognition-mode schemas}
\label{sec:predictions:theorem-grade-lifting}

The schemas of \Cref{sec:recognition} can become theorems only after
their primitive predicates are defined; until then they are neither
provable nor refutable.  Given definitions, each schema needs its own
argument.  \Cref{thm:template-applicability} needs a proof that a
successful seven-stage audit gives the closure the status of a
conditional theorem under its premises.  \Cref{thm:verdict-signature}
needs a derivation of the verdict pattern \SigPattern\ from column
membership.  \Cref{thm:closure-content-thesis} needs, for the
load-bearing content of a trace-state-only closure, a proof that it is
formed-layer content and is named as a source, and a non-derivability
proof showing that the framework's primitives do not yield it.
\Cref{thm:attack-foreclosure-v5} needs, beyond
\Cref{thm:attack-foreclosure-v4}, a proof that every recognition-mode
landing imports the load-bearing content and lands at theorem grade.
Further examples, such as a third trace-state-only application beyond
the P-versus-NP and Riemann-hypothesis cases, would test these
arguments but cannot replace them.  This is the main open problem left
by \Cref{sec:recognition}.

\subsection{A bound on the depth of matrix-element towers}
\label{sec:predictions:ctmt-depth-bound}

\Cref{thm:ctmt-recursion} describes towers of matrix-element
terminations of any finite depth but gives no bound on the depth.  The
open problem is to find an invariant of the carrier
that bounds the depth independently of the application, for example the
complexity of the inherited records, a descent height, the rank of the
native matrix family, or a refinement of the adequacy residual.  Such a
bound cannot simply forbid further reductions by definition; it has to
explain why a deeper matrix-element family cannot arise in a typed
context.

\subsection{Open functorial gates}
\label{sec:predictions:functorial-gates}

The three gates of \Cref{cat:named-functorial-gates},
\(\GateFunctionDegree\), \(\GatePlantedStatistical\) and
\(\GateUniformTime\), each name a transport that has not been proved: from the degree of a
detector to the degree of a Nullstellensatz certificate, from planted
statistical detection to worst-case hardness of a language, and from a
nonuniform family to a uniform time-hierarchy statement.  Each is a
problem in its own right.  The second transports content from the
average case to the worst case, so the corresponding hardness statement
is a worst-case to average-case reduction, and it is related to known
obstacles: a
non-adaptive reduction from an NP-complete problem in the worst case to
a distributional NP problem under a polynomial-time samplable
distribution, tolerant to a small fraction of oracle errors, would
imply \(\mathrm{coNP}\subseteq\mathrm{NP/poly}\)
\cite{BogdanovTrevisan2006WorstAverage}, extending the earlier result
of \cite{FeigenbaumFortnow1993RandomSelfReducibility}; whether a
planted detection problem falls under these hypotheses depends on its
formulation.  Further gates should appear whenever another
application needs a transport between calibrations, regimes, or a
native type and a target type that is not a change of basis, for
instance when its coverage table has a hole like those of
\Cref{thm:coverage-family-geometry}.  The list of gates is not claimed
to be complete.
